\begin{table}[H]
\centering
\caption{Historia de usuario 1: Mostrar menú principal (elaboración propia).\label{tab:hu1}}
\renewcommand{\arraystretch}{1.5}
\begin{tabular}{|p{0.45\textwidth}|p{0.45\textwidth}|}
\hline
\multicolumn{2}{|l|}{\textbf{Historia de Usuario}} \\
\hline
\textbf{Número:} 1 & \textbf{Nombre:} Mostrar menú principal \\
\hline
\textbf{Prioridad:} Alta & \textbf{Riesgo:} Baja \\
\hline
\textbf{Puntos de estimación:} 1.0 & \textbf{Iteración asignada:} 1 \\
\hline
\multicolumn{2}{|p{0.945\textwidth}|}{\textbf{Descripción:} Como usuario quiero acceder a un menú principal intuitivo para navegar por las opciones principales del juego.} \\
\hline
\multicolumn{2}{|p{0.945\textwidth}|}{\textbf{Observaciones:}
\vspace{-0.5em}
\begin{itemize}
    \setlength\itemsep{0em}
    \item El menú principal debe incluir botones: Iniciar Partida, Instrucciones y Salir.
\end{itemize}
\vspace{-1em}
} \\
\hline
\end{tabular}
\end{table}

\begin{table}[H]
\centering
\caption{Historia de usuario 2: Mostrar pantalla de instrucciones (elaboración propia).\label{tab:hu2}}
\renewcommand{\arraystretch}{1.5}
\begin{tabular}{|p{0.45\textwidth}|p{0.45\textwidth}|}
\hline
\multicolumn{2}{|l|}{\textbf{Historia de Usuario}} \\
\hline
\textbf{Número:} 2 & \textbf{Nombre:} Mostrar pantalla de instrucciones \\
\hline
\textbf{Prioridad:} Media & \textbf{Riesgo:} Baja \\
\hline
\textbf{Puntos de estimación:} 1.0 & \textbf{Iteración asignada:} 1 \\
\hline
\multicolumn{2}{|p{0.945\textwidth}|}{\textbf{Descripción:} Como jugador quiero leer las instrucciones del juego para aprender a jugarlo correctamente.} \\
\hline
\multicolumn{2}{|p{0.945\textwidth}|}{\textbf{Observaciones:}
\vspace{-0.5em}
\begin{itemize}
    \setlength\itemsep{0em}
    \item La pantalla debe explicar las reglas de avance, escalera, serpiente y preguntas ODS de forma visual.
    \item Debe ser accesible desde el menú principal.
\end{itemize}
\vspace{-1em}
} \\
\hline
\end{tabular}
\end{table}

\begin{table}[H]
\centering
\caption{Historia de usuario 3: Iniciar partida (elaboración propia).\label{tab:hu3}}
\renewcommand{\arraystretch}{1.5}
\begin{tabular}{|p{0.45\textwidth}|p{0.45\textwidth}|}
\hline
\multicolumn{2}{|l|}{\textbf{Historia de Usuario}} \\
\hline
\textbf{Número:} 3 & \textbf{Nombre:} Iniciar partida \\
\hline
\textbf{Prioridad:} Alta & \textbf{Riesgo:} Medio \\
\hline
\textbf{Puntos de estimación:} 2.0 & \textbf{Iteración asignada:} 1 \\
\hline
\multicolumn{2}{|p{0.945\textwidth}|}{\textbf{Descripción:} Como jugador quiero iniciar una nueva partida para aprender sobre los ODS de manera interactiva.} \\
\hline
\multicolumn{2}{|p{0.945\textwidth}|}{\textbf{Observaciones:}
\vspace{-0.5em}
\begin{itemize}
    \setlength\itemsep{0em}
    \item Debe incluir botón para iniciar partida.
    \item Permite configurar número de jugadores (1--4).
    \item Carga el tablero y realiza transición a la primera ronda.
    \item Se inicializan estadísticas en cero.
\end{itemize}
\vspace{-1em}
} \\
\hline
\end{tabular}
\end{table}

\begin{table}[H]
\centering
\caption{Historia de usuario 4: Visualizar tablero (elaboración propia).\label{tab:hu4}}
\renewcommand{\arraystretch}{1.5}
\begin{tabular}{|p{0.45\textwidth}|p{0.45\textwidth}|}
\hline
\multicolumn{2}{|l|}{\textbf{Historia de Usuario}} \\
\hline
\textbf{Número:} 4 & \textbf{Nombre:} Visualizar tablero \\
\hline
\textbf{Prioridad:} Alta & \textbf{Riesgo:} Alto \\
\hline
\textbf{Puntos de estimación:} 5.0 & \textbf{Iteración asignada:} 1 \\
\hline
\multicolumn{2}{|p{0.945\textwidth}|}{\textbf{Descripción:} Como jugador quiero ver el tablero completo de 63 casillas para entender el recorrido hacia el año 2030.} \\
\hline
\multicolumn{2}{|p{0.945\textwidth}|}{\textbf{Observaciones:}
\vspace{-0.5em}
\begin{itemize}
    \setlength\itemsep{0em}
    \item El tablero debe ser fiel al diseño oficial.
    \item Las 17 casillas temáticas deben identificarse con el ícono del ODS correspondiente.
    \item El tablero debe ser visible sin \textit{scroll}.
    \item Debe ser adaptable a distintas resoluciones de pantalla.
\end{itemize}
\vspace{-1em}
} \\
\hline
\end{tabular}
\end{table}

\begin{table}[H]
\centering
\caption{Historia de usuario 5: Lanzar dado (elaboración propia).\label{tab:hu5}}
\renewcommand{\arraystretch}{1.5}
\begin{tabular}{|p{0.45\textwidth}|p{0.45\textwidth}|}
\hline
\multicolumn{2}{|l|}{\textbf{Historia de Usuario}} \\
\hline
\textbf{Número:} 5 & \textbf{Nombre:} Lanzar dado \\
\hline
\textbf{Prioridad:} Alta & \textbf{Riesgo:} Medio \\
\hline
\textbf{Puntos de estimación:} 1.5 & \textbf{Iteración asignada:} 2 \\
\hline
\multicolumn{2}{|p{0.945\textwidth}|}{\textbf{Descripción:} Como jugador quiero lanzar un dado virtual para determinar el avance de mi ficha en el tablero.} \\
\hline
\multicolumn{2}{|p{0.945\textwidth}|}{\textbf{Observaciones:}
\vspace{-0.5em}
\begin{itemize}
    \setlength\itemsep{0em}
    \item Debe mostrar animación breve del dado.
    \item El resultado entre 1 y 6 debe ser pseudoaleatorio.
    \item No debe ser posible lanzar el dado más de una vez por turno.
    \item Debe incluir \textit{feedback} visual y sonoro al obtener el resultado.
\end{itemize}
\vspace{-1em}
} \\
\hline
\end{tabular}
\end{table}

\begin{table}[H]
\centering
\caption{Historia de usuario 6: Mover ficha (elaboración propia).\label{tab:hu6}}
\renewcommand{\arraystretch}{1.5}
\begin{tabular}{|p{0.45\textwidth}|p{0.45\textwidth}|}
\hline
\multicolumn{2}{|l|}{\textbf{Historia de Usuario}} \\
\hline
\textbf{Número:} 6 & \textbf{Nombre:} Mover ficha \\
\hline
\textbf{Prioridad:} Alta & \textbf{Riesgo:} Alto \\
\hline
\textbf{Puntos de estimación:} 4.0 & \textbf{Iteración asignada:} 2 \\
\hline
\multicolumn{2}{|p{0.945\textwidth}|}{\textbf{Descripción:} Como jugador quiero que mi ficha se mueva automáticamente según el resultado del dado para concentrarme en el contenido educativo.} \\
\hline
\multicolumn{2}{|p{0.945\textwidth}|}{\textbf{Observaciones:}
\vspace{-0.5em}
\begin{itemize}
    \setlength\itemsep{0em}
    \item El movimiento debe ser animado casilla por casilla.
    \item Debe calcularse la trayectoria correcta.
    \item Al llegar a la casilla final se comprueba si es casilla ODS, escalera o serpiente.
\end{itemize}
\vspace{-1em}
} \\
\hline
\end{tabular}
\end{table}

\begin{table}[H]
\centering
\caption{Historia de usuario 7: Aplicar efecto de escalera (elaboración propia).\label{tab:hu7}}
\renewcommand{\arraystretch}{1.5}
\begin{tabular}{|p{0.45\textwidth}|p{0.45\textwidth}|}
\hline
\multicolumn{2}{|l|}{\textbf{Historia de Usuario}} \\
\hline
\textbf{Número:} 7 & \textbf{Nombre:} Aplicar efecto de escalera \\
\hline
\textbf{Prioridad:} Alta & \textbf{Riesgo:} Medio \\
\hline
\textbf{Puntos de estimación:} 1.5 & \textbf{Iteración asignada:} 2 \\
\hline
\multicolumn{2}{|p{0.945\textwidth}|}{\textbf{Descripción:} Como jugador quiero que mi ficha avance automáticamente por una escalera al caer en su base para ser recompensado por mi progreso.} \\
\hline
\multicolumn{2}{|p{0.945\textwidth}|}{\textbf{Observaciones:}
\vspace{-0.5em}
\begin{itemize}
    \setlength\itemsep{0em}
    \item El sistema debe detectar si la casilla final es la base de una escalera.
    \item La ficha debe avanzar a la casilla destino con animación de subida.
    \item Debe incluir \textit{feedback} sonoro positivo.
\end{itemize}
\vspace{-1em}
} \\
\hline
\end{tabular}
\end{table}

\begin{table}[H]
\centering
\caption{Historia de usuario 8: Aplicar efecto de serpiente (elaboración propia).\label{tab:hu8}}
\renewcommand{\arraystretch}{1.5}
\begin{tabular}{|p{0.45\textwidth}|p{0.45\textwidth}|}
\hline
\multicolumn{2}{|l|}{\textbf{Historia de Usuario}} \\
\hline
\textbf{Número:} 8 & \textbf{Nombre:} Aplicar efecto de serpiente \\
\hline
\textbf{Prioridad:} Alta & \textbf{Riesgo:} Medio \\
\hline
\textbf{Puntos de estimación:} 1.5 & \textbf{Iteración asignada:} 2 \\
\hline
\multicolumn{2}{|p{0.945\textwidth}|}{\textbf{Descripción:} Como jugador quiero que mi ficha retroceda automáticamente al caer en la cabeza de una serpiente para aplicar la penalización del juego clásico.} \\
\hline
\multicolumn{2}{|p{0.945\textwidth}|}{\textbf{Observaciones:}
\vspace{-0.5em}
\begin{itemize}
    \setlength\itemsep{0em}
    \item El sistema debe detectar si la casilla final es la cabeza de una serpiente.
    \item La ficha retrocede a la casilla destino con animación descendente.
    \item Debe incluir \textit{feedback} sonoro negativo.
\end{itemize}
\vspace{-1em}
} \\
\hline
\end{tabular}
\end{table}

\begin{table}[H]
\centering
\caption{Historia de usuario 10: Evaluar respuesta (elaboración propia).\label{tab:hu10}}
\renewcommand{\arraystretch}{1.5}
\begin{tabular}{|p{0.45\textwidth}|p{0.45\textwidth}|}
\hline
\multicolumn{2}{|l|}{\textbf{Historia de Usuario}} \\
\hline
\textbf{Número:} 10 & \textbf{Nombre:} Evaluar respuesta \\
\hline
\textbf{Prioridad:} Alta & \textbf{Riesgo:} Medio \\
\hline
\textbf{Puntos de estimación:} 3.0 & \textbf{Iteración asignada:} 3 \\
\hline
\multicolumn{2}{|p{0.945\textwidth}|}{\textbf{Descripción:} Como jugador quiero recibir retroalimentación inmediata al responder para entender si acerté y por qué.} \\
\hline
\multicolumn{2}{|p{0.945\textwidth}|}{\textbf{Observaciones:}
\vspace{-0.5em}
\begin{itemize}
    \setlength\itemsep{0em}
    \item El sistema debe mostrar un indicador visual de resultado (verde/rojo).
    \item Debe mostrar un mensaje explicativo breve de la respuesta correcta.
    \item Si la respuesta es correcta, el jugador recibe una bonificación obteniendo la oportunidad de volver a lanzar los dados. Si es incorrecta, el turno finaliza y pasa al siguiente jugador.
\end{itemize}
\vspace{-1em}
} \\
\hline
\end{tabular}
\end{table}

\begin{table}[H]
\centering
\caption{Historia de usuario 11: Gestionar turnos (elaboración propia).\label{tab:hu11}}
\renewcommand{\arraystretch}{1.5}
\begin{tabular}{|p{0.45\textwidth}|p{0.45\textwidth}|}
\hline
\multicolumn{2}{|l|}{\textbf{Historia de Usuario}} \\
\hline
\textbf{Número:} 11 & \textbf{Nombre:} Gestionar turnos \\
\hline
\textbf{Prioridad:} Alta & \textbf{Riesgo:} Baja \\
\hline
\textbf{Puntos de estimación:} 2.0 & \textbf{Iteración asignada:} 3 \\
\hline
\multicolumn{2}{|p{0.945\textwidth}|}{\textbf{Descripción:} Como jugador quiero que el sistema gestione turnos automáticamente y muestre quién es el jugador activo para que el juego fluya.} \\
\hline
\multicolumn{2}{|p{0.945\textwidth}|}{\textbf{Observaciones:}
\vspace{-0.5em}
\begin{itemize}
    \setlength\itemsep{0em}
    \item El sistema debe indicar visualmente el jugador activo resaltando su ficha.
    \item Al finalizar el turno debe pasar automáticamente al siguiente jugador en orden.
\end{itemize}
\vspace{-1em}
} \\
\hline
\end{tabular}
\end{table}

\begin{table}[H]
\centering
\caption{Historia de usuario 12: Determinar condición de victoria (elaboración propia).\label{tab:hu12}}
\renewcommand{\arraystretch}{1.5}
\begin{tabular}{|p{0.45\textwidth}|p{0.45\textwidth}|}
\hline
\multicolumn{2}{|l|}{\textbf{Historia de Usuario}} \\
\hline
\textbf{Número:} 12 & \textbf{Nombre:} Determinar condición de victoria \\
\hline
\textbf{Prioridad:} Alta & \textbf{Riesgo:} Baja \\
\hline
\textbf{Puntos de estimación:} 1.5 & \textbf{Iteración asignada:} 4 \\
\hline
\multicolumn{2}{|p{0.945\textwidth}|}{\textbf{Descripción:} Como jugador quiero que el juego detecte cuando alguien llegue a la casilla 63 para mostrar la pantalla de victoria.} \\
\hline
\multicolumn{2}{|p{0.945\textwidth}|}{\textbf{Observaciones:}
\vspace{-0.5em}
\begin{itemize}
    \setlength\itemsep{0em}
    \item El sistema debe detectar la llegada exacta a la casilla final.
    \item Debe mostrar una pantalla de victoria con animación de confeti y el nombre del ganador.
\end{itemize}
\vspace{-1em}
} \\
\hline
\end{tabular}
\end{table}

\begin{table}[H]
\centering
\caption{Historia de usuario 13: Mostrar estadísticas de partida (elaboración propia).\label{tab:hu13}}
\renewcommand{\arraystretch}{1.5}
\begin{tabular}{|p{0.45\textwidth}|p{0.45\textwidth}|}
\hline
\multicolumn{2}{|l|}{\textbf{Historia de Usuario}} \\
\hline
\textbf{Número:} 13 & \textbf{Nombre:} Mostrar estadísticas de partida \\
\hline
\textbf{Prioridad:} Media & \textbf{Riesgo:} Medio \\
\hline
\textbf{Puntos de estimación:} 4.0 & \textbf{Iteración asignada:} 4 \\
\hline
\multicolumn{2}{|p{0.945\textwidth}|}{\textbf{Descripción:} Como educador quiero visualizar estadísticas de desempeño por jugador y por ODS de manera continua en tiempo real para evaluar el aprendizaje progresivamente.} \\
\hline
\multicolumn{2}{|p{0.945\textwidth}|}{\textbf{Observaciones:}
\vspace{-0.5em}
\begin{itemize}
    \setlength\itemsep{0em}
    \item Debe existir una interfaz accesible que muestre las estadísticas de desempeño de cada jugador en tiempo real.
    \item Debe incluir el porcentaje de respuestas correctas por jugador.
    \item Debe identificar los ODS con menor porcentaje de aciertos para cada jugador.
\end{itemize}
\vspace{-1em}
} \\
\hline
\end{tabular}
\end{table}

\begin{table}[H]
\centering
\caption{Historia de usuario 14: Configurar audio (elaboración propia).\label{tab:hu14}}
\renewcommand{\arraystretch}{1.5}
\begin{tabular}{|p{0.45\textwidth}|p{0.45\textwidth}|}
\hline
\multicolumn{2}{|l|}{\textbf{Historia de Usuario}} \\
\hline
\textbf{Número:} 14 & \textbf{Nombre:} Configurar audio \\
\hline
\textbf{Prioridad:} Baja & \textbf{Riesgo:} Baja \\
\hline
\textbf{Puntos de estimación:} 2.5 & \textbf{Iteración asignada:} 4 \\
\hline
\multicolumn{2}{|p{0.945\textwidth}|}{\textbf{Descripción:} Como jugador quiero poder activar o desactivar la música y los efectos de sonido para personalizar mi partida.} \\
\hline
\multicolumn{2}{|p{0.945\textwidth}|}{\textbf{Observaciones:}
\vspace{-0.5em}
\begin{itemize}
    \setlength\itemsep{0em}
    \item Debe incluir controles para el volumen global, música y efectos sonoros.
\end{itemize}
\vspace{-1em}
} \\
\hline
\end{tabular}
\end{table}

\begin{table}[H]
\centering
\caption{Historia de usuario 15: Pausar partida (elaboración propia).\label{tab:hu15}}
\renewcommand{\arraystretch}{1.5}
\begin{tabular}{|p{0.45\textwidth}|p{0.45\textwidth}|}
\hline
\multicolumn{2}{|l|}{\textbf{Historia de Usuario}} \\
\hline
\textbf{Número:} 15 & \textbf{Nombre:} Pausar partida \\
\hline
\textbf{Prioridad:} Media & \textbf{Riesgo:} Baja \\
\hline
\textbf{Puntos de estimación:} 1.0 & \textbf{Iteración asignada:} 4 \\
\hline
\multicolumn{2}{|p{0.945\textwidth}|}{\textbf{Descripción:} Como jugador quiero pausar la partida en cualquier momento para atender interrupciones sin perder el progreso.} \\
\hline
\multicolumn{2}{|p{0.945\textwidth}|}{\textbf{Observaciones:}
\vspace{-0.5em}
\begin{itemize}
    \setlength\itemsep{0em}
    \item Debe existir un panel de pausa que detenga todo el flujo y las animaciones.
    \item Debe permitir reanudar desde el estado exacto en que se pausó.
\end{itemize}
\vspace{-1em}
} \\
\hline
\end{tabular}
\end{table}

\begin{table}[H]
\centering
\caption{Historia de usuario 16: Reiniciar partida (elaboración propia).\label{tab:hu16}}
\renewcommand{\arraystretch}{1.5}
\begin{tabular}{|p{0.45\textwidth}|p{0.45\textwidth}|}
\hline
\multicolumn{2}{|l|}{\textbf{Historia de Usuario}} \\
\hline
\textbf{Número:} 16 & \textbf{Nombre:} Reiniciar partida \\
\hline
\textbf{Prioridad:} Media & \textbf{Riesgo:} Baja \\
\hline
\textbf{Puntos de estimación:} 1.0 & \textbf{Iteración asignada:} 4 \\
\hline
\multicolumn{2}{|p{0.945\textwidth}|}{\textbf{Descripción:} Como jugador quiero reiniciar la partida actual desde cero sin tener que volver a configurar el número de jugadores.} \\
\hline
\multicolumn{2}{|p{0.945\textwidth}|}{\textbf{Observaciones:}
\vspace{-0.5em}
\begin{itemize}
    \setlength\itemsep{0em}
    \item Accesible desde el menú de pausa o de victoria.
    \item Debe restaurar las posiciones y vaciar las estadísticas de la sesión actual.
\end{itemize}
\vspace{-1em}
} \\
\hline
\end{tabular}
\end{table}

\begin{table}[H]
\centering
\caption{Historia de usuario 17: Habilitar cámara de seguimiento (elaboración propia).\label{tab:hu17}}
\renewcommand{\arraystretch}{1.5}
\begin{tabular}{|p{0.45\textwidth}|p{0.45\textwidth}|}
\hline
\multicolumn{2}{|l|}{\textbf{Historia de Usuario}} \\
\hline
\textbf{Número:} 17 & \textbf{Nombre:} Habilitar cámara de seguimiento \\
\hline
\textbf{Prioridad:} Alta & \textbf{Riesgo:} Medio \\
\hline
\textbf{Puntos de estimación:} 2.0 & \textbf{Iteración asignada:} 5 \\
\hline
\multicolumn{2}{|p{0.945\textwidth}|}{\textbf{Descripción:} Como jugador quiero que la cámara enfoque automáticamente a mi ficha cuando sea mi turno para tener mejor inmersión y saber que me toca.} \\
\hline
\multicolumn{2}{|p{0.945\textwidth}|}{\textbf{Observaciones:}
\vspace{-0.5em}
\begin{itemize}
    \setlength\itemsep{0em}
    \item La cámara debe realizar un \textit{zoom} suave y centrarse en el jugador activo al inicio de cada turno.
    \item Durante el movimiento de la ficha, la cámara debe seguir su trayectoria en tiempo real.
    \item Al finalizar el turno de un jugador, la cámara debe transicionar automáticamente hacia la ubicación de la ficha del siguiente jugador activo.
\end{itemize}
\vspace{-1em}
} \\
\hline
\end{tabular}
\end{table}

\begin{table}[H]
\centering
\caption{Historia de usuario 18: Habilitar modo de práctica (elaboración propia).\label{tab:hu18}}
\renewcommand{\arraystretch}{1.5}
\begin{tabular}{|p{0.45\textwidth}|p{0.45\textwidth}|}
\hline
\multicolumn{2}{|l|}{\textbf{Historia de Usuario}} \\
\hline
\textbf{Número:} 18 & \textbf{Nombre:} Habilitar modo de práctica \\
\hline
\textbf{Prioridad:} Media & \textbf{Riesgo:} Bajo \\
\hline
\textbf{Puntos de estimación:} 5.0 & \textbf{Iteración asignada:} 5 \\
\hline
\multicolumn{2}{|p{0.945\textwidth}|}{\textbf{Descripción:} Como usuario nuevo quiero poder jugar la partida completamente solo, recorriendo el tablero a mi propio ritmo para practicar y conocer mis habilidades.} \\
\hline
\multicolumn{2}{|p{0.945\textwidth}|}{\textbf{Observaciones:}
\vspace{-0.5em}
\begin{itemize}
    \setlength\itemsep{0em}
    \item El entorno conserva todo el ciclo de progreso, escaleras y rondas de preguntas educativas pero restringido y centrado únicamente en el desarrollo individual del jugador.
\end{itemize}
\vspace{-1em}
} \\
\hline
\end{tabular}
\end{table}
